\section*{M24/M25: Sound finite-batch certificates by restricted gains and Gram refinement}
\label{sec:m24-m25-finite-batch-certificates}

\paragraph{Scope.}
These are \emph{finite-batch} statements. Everything below is quantified over
a fixed finite family of leaf data and a fixed rank assignment; nothing here
is a claim about unbounded inputs, and neither certificate is a computational
shortcut (both evaluations are needed to form the bound). Their content is
that the root discrepancy is \emph{compositionally attributable} to per-node
sources with explicit transport factors, and that two specific composition
rules are sound and monotonically ordered.

\paragraph{Setup.}
Let $T$ be a finite ordered rooted tree with root $r$, internal vertices
carrying multilinear laws $\mu_v$ and orthogonal projectors $P_v$ of
prescribed rank. Fix a finite index set $n\in\{1,\dots,N\}$ of samples, and
for each $n$ leaf data $z_\ell^{(n)}$. Define, for each $n$:
\[
  F_\ell^{(n)}:=z_\ell^{(n)},\qquad
  F_v^{(n)}:=\mu_v\bigl(F_{c_1(v)}^{(n)},\dots,F_{c_m(v)}^{(n)}\bigr),
\]
\[
  R_\ell^{(n)}:=z_\ell^{(n)},\qquad
  \widetilde R_v^{(n)}:=\mu_v\bigl(R_{c_1(v)}^{(n)},\dots,R_{c_m(v)}^{(n)}\bigr),
  \qquad
  R_v^{(n)}:=
  \begin{cases}
    P_v\widetilde R_v^{(n)}, & v\ne r,\\
    \widetilde R_v^{(n)}, & v=r.
  \end{cases}
\]
The root is not projected, matching the convention under which the
projected-root discrepancy is measured in the root's own space. Put
\[
  \delta_v^{(n)}:=F_v^{(n)}-R_v^{(n)}\quad(v\ne r),
  \qquad
  \delta_\ell^{(n)}:=0,
  \qquad
  E^{(n)}:=F_r^{(n)}-\widetilde R_r^{(n)}.
\]

\paragraph{Definition (effective slot map).}
For an internal $v$ with children $c_1,\dots,c_m$, a slot $j$, and a sample
$n$, let $M_{v,j}^{(n)}:H_{c_j}\to H_v$ be the linear map
\[
  \boxed{\;
  M_{v,j}^{(n)}(u)
  :=\mu_v\bigl(R_{c_1}^{(n)},\dots,R_{c_{j-1}}^{(n)},\,u,\,
               F_{c_{j+1}}^{(n)},\dots,F_{c_m}^{(n)}\bigr).\;}
\tag{1}
\]
Slots below $j$ carry reduced values, slots above $j$ carry ambient values.
Linearity in $u$ is immediate from multilinearity of $\mu_v$.

\paragraph{Lemma 1 (exact telescoping).}
For every internal $v$ and every $n$,
\[
  F_v^{(n)}-\widetilde R_v^{(n)}
  =\sum_{j=1}^{m}M_{v,j}^{(n)}\bigl(\delta_{c_j}^{(n)}\bigr).
\tag{2}
\]

\emph{Proof.} Write $x_i:=F_{c_i}^{(n)}$, $y_i:=R_{c_i}^{(n)}$. Then
\[
 \mu_v(x_1,\dots,x_m)-\mu_v(y_1,\dots,y_m)
 =\sum_{j=1}^m\Bigl[\mu_v(y_{<j},x_j,x_{>j})-\mu_v(y_{<j},y_j,x_{>j})\Bigr],
\]
a telescoping sum in which consecutive terms cancel. By multilinearity in
slot $j$ the $j$-th bracket equals $\mu_v(y_{<j},x_j-y_j,x_{>j})
=M_{v,j}^{(n)}(x_j-y_j)$, and $x_j-y_j=\delta_{c_j}^{(n)}$. $\square$

Note that (2) is an identity, not an estimate: no slack has been introduced
yet. All subsequent looseness is attributable to how (2) is bounded.

\paragraph{Definition (M24 bound).}
Set $B_\ell^{(n)}:=0$ at leaves, and recursively in postorder
\[
  \gamma_v^{(n)}:=\bigl\|(I-P_v)\widetilde R_v^{(n)}\bigr\|,
  \qquad
  p_v^{(n)}:=\sum_{j=1}^m\bigl\|M_{v,j}^{(n)}\bigr\|_{\mathrm{op}}\,B_{c_j}^{(n)},
\]
\[
  B_v^{(n)}:=
  \begin{cases}
    p_v^{(n)}+\gamma_v^{(n)}, & v\ne r,\\
    p_r^{(n)}, & v=r.
  \end{cases}
\tag{3}
\]

\paragraph{Theorem M24 (per-sample soundness).}
For every $n$ and every internal $v\ne r$, $\;\|\delta_v^{(n)}\|\le B_v^{(n)}$,
and at the root
\[
  \boxed{\;\bigl\|E^{(n)}\bigr\|\le B_r^{(n)}.\;}
\tag{4}
\]

\emph{Proof.} Induction on the postorder. At a leaf both sides vanish. Let
$v$ be internal and assume the claim at every child. By Lemma 1 and
subadditivity,
\[
 \bigl\|F_v^{(n)}-\widetilde R_v^{(n)}\bigr\|
 \le\sum_j\bigl\|M_{v,j}^{(n)}\bigr\|_{\mathrm{op}}\bigl\|\delta_{c_j}^{(n)}\bigr\|
 \le\sum_j\bigl\|M_{v,j}^{(n)}\bigr\|_{\mathrm{op}}B_{c_j}^{(n)}=p_v^{(n)},
\tag{5}
\]
the second step using the inductive hypothesis and nonnegativity of operator
norms. If $v=r$ then $E^{(n)}=F_r^{(n)}-\widetilde R_r^{(n)}$ and (4) follows
from (5). If $v\ne r$, decompose
\[
 \delta_v^{(n)}
 =F_v^{(n)}-P_v\widetilde R_v^{(n)}
 =\underbrace{\bigl(F_v^{(n)}-\widetilde R_v^{(n)}\bigr)}_{=:a_v^{(n)}}
 +\underbrace{(I-P_v)\widetilde R_v^{(n)}}_{=:b_v^{(n)}},
\tag{6}
\]
so $\|\delta_v^{(n)}\|\le\|a_v^{(n)}\|+\|b_v^{(n)}\|\le p_v^{(n)}+\gamma_v^{(n)}$.
$\square$

\paragraph{Corollary M24.1 (sound reduction).}
$\displaystyle\max_n\bigl\|E^{(n)}\bigr\|\le\max_n B_r^{(n)}$.
Immediate from (4).

\paragraph{Proposition M24.2 (per-sample propagation is never looser).}
Let $\bar B$ denote the batch-decoupled variant, obtained by maximizing each
ingredient over the batch before composing:
\[
  \bar B_\ell:=0,\qquad
  \bar B_v:=\sum_j\Bigl(\max_n\bigl\|M_{v,j}^{(n)}\bigr\|_{\mathrm{op}}\Bigr)\bar B_{c_j}
  +\max_n\gamma_v^{(n)}\quad(v\ne r).
\]
Then $\max_nB_v^{(n)}\le\bar B_v$ for every $v$.

\emph{Proof.} Induction. For nonnegative reals,
$\max_n\sum_j\alpha_j^{(n)}\beta_j^{(n)}\le\sum_j(\max_n\alpha_j^{(n)})(\max_n\beta_j^{(n)})$,
since each term is bounded by the corresponding product of maxima and the
maximum of a sum is at most the sum of maxima. Applying this to (3) and using
$\max_nB_{c_j}^{(n)}\le\bar B_{c_j}$ gives the claim. $\square$

The inequality is generally strict: $\bar B$ pairs the sample maximizing a
gain with the (different) sample maximizing the incoming error, a
configuration realized by no single sample. Since the defect is incurred once
per level it compounds multiplicatively in the depth.

\paragraph{Definition (M25 bound).}
Retain $\gamma_v^{(n)}$ and define, with $G_\ell^{(n)}:=0$,
\[
  \pi_v^{(n)}:=\sum_j\bigl\|M_{v,j}^{(n)}\bigr\|_{\mathrm{op}}G_{c_j}^{(n)},
  \qquad
  \nu_v^{(n)}:=\sum_j\bigl\|(I-P_v)M_{v,j}^{(n)}\bigr\|_{\mathrm{op}}G_{c_j}^{(n)},
\]
\[
  \kappa_v^{(n)}:=\min\bigl(\nu_v^{(n)},\pi_v^{(n)}\bigr),
  \qquad
  G_v^{(n)}:=
  \begin{cases}
   \sqrt{\bigl(\pi_v^{(n)}\bigr)^2+\bigl(\gamma_v^{(n)}\bigr)^2
         +2\kappa_v^{(n)}\gamma_v^{(n)}}, & v\ne r,\\[1mm]
   \pi_r^{(n)}, & v=r.
  \end{cases}
\tag{7}
\]

\paragraph{Theorem M25 (Gram refinement).}
For every $n$ and every $v$,
\[
  \boxed{\;\bigl\|\delta_v^{(n)}\bigr\|\le G_v^{(n)}\le B_v^{(n)}\;}
\tag{8}
\]
(with $\|E^{(n)}\|\le G_r^{(n)}$ at the root), where $B$ is computed by (3).

\emph{Proof.} Soundness, by induction. The bound (5) holds verbatim with
$B_{c_j}$ replaced by $G_{c_j}$, giving $\|a_v^{(n)}\|\le\pi_v^{(n)}$; the
root case follows. For $v\ne r$ use the orthogonal decomposition (6). Since
$P_v$ is an orthogonal projector, $I-P_v$ is self-adjoint and idempotent and
$b_v^{(n)}=(I-P_v)b_v^{(n)}$, whence
\[
 \bigl\langle a_v^{(n)},b_v^{(n)}\bigr\rangle
 =\bigl\langle (I-P_v)a_v^{(n)},b_v^{(n)}\bigr\rangle
 \le\bigl\|(I-P_v)a_v^{(n)}\bigr\|\,\bigl\|b_v^{(n)}\bigr\|.
\]
By Lemma 1, $(I-P_v)a_v^{(n)}=\sum_j(I-P_v)M_{v,j}^{(n)}(\delta_{c_j}^{(n)})$,
so $\|(I-P_v)a_v^{(n)}\|\le\nu_v^{(n)}$. Combining with the Cauchy--Schwarz
bound $\langle a_v,b_v\rangle\le\|a_v\|\|b_v\|\le\pi_v^{(n)}\gamma_v^{(n)}$
gives $\langle a_v,b_v\rangle\le\kappa_v^{(n)}\gamma_v^{(n)}$. Therefore
\[
 \bigl\|\delta_v^{(n)}\bigr\|^2
 =\bigl\|a_v^{(n)}\bigr\|^2+\bigl\|b_v^{(n)}\bigr\|^2
  +2\bigl\langle a_v^{(n)},b_v^{(n)}\bigr\rangle
 \le\bigl(\pi_v^{(n)}\bigr)^2+\bigl(\gamma_v^{(n)}\bigr)^2
  +2\kappa_v^{(n)}\gamma_v^{(n)},
\]
which is $\bigl(G_v^{(n)}\bigr)^2$.

Ordering: $\kappa_v^{(n)}\le\pi_v^{(n)}$ by definition, so
\[
 \bigl(G_v^{(n)}\bigr)^2
 \le\bigl(\pi_v^{(n)}\bigr)^2+\bigl(\gamma_v^{(n)}\bigr)^2
   +2\pi_v^{(n)}\gamma_v^{(n)}
 =\bigl(\pi_v^{(n)}+\gamma_v^{(n)}\bigr)^2 ,
\]
and $\pi_v^{(n)}\le p_v^{(n)}$ follows inductively from $G_{c_j}\le B_{c_j}$.
$\square$

\paragraph{Proposition M25.1 (equality and strict improvement).}
At a node $v\ne r$ the refinement is strict, $G_v^{(n)}<\pi_v^{(n)}+\gamma_v^{(n)}$,
if and only if
\[
  \gamma_v^{(n)}>0
  \quad\text{and}\quad
  \nu_v^{(n)}<\pi_v^{(n)} .
\]
Since $\|(I-P_v)M\|_{\mathrm{op}}\le\|M\|_{\mathrm{op}}$ always, the
refinement is never worse. Equality $\nu_v^{(n)}=\pi_v^{(n)}$ holds exactly
when, for each $j$ contributing a nonzero $G_{c_j}^{(n)}$, some maximizing
singular direction of $M_{v,j}^{(n)}$ has image entirely inside
$\operatorname{Ran}(I-P_v)$.

\emph{Magnitude.} When $\nu_v\ll\pi_v$ the bound degrades to
$\sqrt{\pi_v^2+\gamma_v^2}$, i.e.\ Pythagorean rather than linear addition;
the gain over $\pi_v+\gamma_v$ is then largest when $\pi_v\approx\gamma_v$
(a factor $1/\sqrt2$) and vanishes when either term dominates. This is the
finite-batch counterpart of the $k=3$ orthogonality mechanism
$D_1\perp R_1$, and, as there, it is a second-order effect: it does not
change the leading accumulation rate.

\paragraph{Proposition M24.3 (dominance over the Frobenius certificate).}
Let $A$ denote the scalar certificate that propagates value bounds
$U_\ell:=\max_n\|z_\ell^{(n)}\|$, $U_v:=\|\mu_v\|_F\prod_iU_{c_i}$, uses slot
gains $\|\mu_v\|_F\prod_{i\ne j}U_{c_i}$ and closures $\max_n\gamma_v^{(n)}$,
combined additively. Then $B_v^{(n)}\le A_v$ for every $v,n$; hence
$\max_nG_r^{(n)}\le\max_nB_r^{(n)}\le A_r$.

\emph{Proof.} First, $U_v\ge\max_n\max\bigl(\|F_v^{(n)}\|,\|R_v^{(n)}\|\bigr)$
by induction, using $\|\mu_v(x_1,\dots,x_m)\|\le\|\mu_v\|_F\prod_i\|x_i\|$
(the Frobenius norm of the unfolding dominates the multilinear operator norm)
and $\|P_v\|_{\mathrm{op}}\le1$. Consequently, from (1),
\[
 \bigl\|M_{v,j}^{(n)}\bigr\|_{\mathrm{op}}
 \le\|\mu_v\|_F\prod_{i\ne j}\bigl\|(\text{slot }i\text{ value})^{(n)}\bigr\|
 \le\|\mu_v\|_F\prod_{i\ne j}U_{c_i},
\]
which is exactly $A$'s slot gain, and $\gamma_v^{(n)}\le\max_n\gamma_v^{(n)}$.
Induction on the postorder then gives $B_v^{(n)}\le A_v$. $\square$

\paragraph{Boundary.}
The certified chain is
$\|E^{(n)}\|\le G_r^{(n)}\le B_r^{(n)}$ and
$\max_nB_r^{(n)}\le\min(\bar B_r,\,A_r)$.
None of these statements bounds $C_T^P(\eta)$ or any other extremal constant:
they are instance-wise, for the given laws, projectors, ranks and batch. In
particular they do not contradict the coherent-accumulation witness for the
free-law class, and they say nothing about how $B_r$ behaves as depth grows
--- empirically it still grows geometrically, at a rate governed by the
step $\|M\delta\|\le\sigma_{\max}(M)\|\delta\|$ in (5), which is the only
remaining source of slack in Lemma 1's otherwise exact expansion.
Replacing that step by directional (subspace or ellipsoid) transport is the
subject of M28 and is not claimed here.
